\begin{textsample}{NEG Dim 9 – global\_north\_2025\_05 – Score -6.00 – global\_north\_2025\_05/125317264.txt}  \label{ex:f9_neg_092}
A raid on the house of British pro-Palestine journalist Asa Winstanley has been ruled unlawful by a British court.
The Central Criminal Court ruled earlier this week that the Metropolitan Police needed to immediately return all devices seized during the October 2024 raid on the writer ' s home, which included computers and phones.
Solicitor Tayab Ali called the ruling a victory for press freedom and accused the police of"silencing a journalist who had made comments on the situation in Gaza"."This ruling is a resounding victory for press freedom and the rule of law,"he said, according to Solicitors Journal."The actions of the police, raiding a journalist ' s home under the guise of counter-terrorism, were not only unlawful but a grave threat to the democratic principle that journalists must be able to work without fear of state harassment."New \textbf{MEE} \textbf{newsletter} : Jerusalem Dispatch Sign up to get the latest insights and analysis on Israel-Palestine, alongside Turkey \textbf{Unpacked} and other \textbf{MEE} \textbf{newsletters} Recorder @ @ @ @ @ @ @ @ @ @ at the Central Criminal Court, stated he was"very troubled by the way in which the search warrant application was drafted, approved and granted where items were to be seized from a journalist."Winstanley, who has written for the \textbf{Electronic} Intifada since 2009 and has been its associate editor since 2012, regularly uses social media to comment on issues relating to Palestine and Israel, including the ongoing war in Gaza. ' This ruling is a resounding victory for press freedom and the rule of law ' Tayab Ali, solicitor Winstanley contributed a number of articles to Middle East Eye between 2015 and 2018.
He is the author of Weaponising Anti-Semitism, which accused the pro-Israel lobby in the UK of orchestrating a campaign to undermine former Labour leader Jeremy Corbyn.
Last year, the general secretaries of the UK ' s National Union of Journalists and the International Federation of Journalists wrote to Metropolitan Police Assistance Commissioner Matt Jukes, the head of Counter Terrorism Policing in the UK, to raise concerns @ @ @ @ @ @ @ @ @ @ The letter came after freelance journalist Richard Medhurst was detained and questioned by police at Heathrow Airport.
Medhurst said he believed he was targeted for speaking out on the situation in Palestine.
Middle East Eye delivers independent and \textbf{unrivalled} coverage and analysis of the Middle East, North Africa and beyond.
To learn more about \textbf{republishing} this content and the associated fees, please fill out this form.
More about \textbf{MEE} can be found here.

% matched lemmas: electronic, mee, newsletter, republish, unpack, unrivalled
\end{textsample}
